%%%%%%%%%%%%%%%%%%%%%%%%%%%%%%%%%%%%%%%%%%%%%%%%%%%%%%%%%%%%%%%%%%%%%%%%%%%%
%% Text Area : 8in (include Runningheads) x 5in
%% ws-jbcb.tex : 18-05-2026
%% Tex file to use with ws-jbcb.cls written in LaTeX2e.
%% The content, structure, format and layout of this style file is the
%% property of World Scientific Publishing Co. Pte. Ltd.
%% Copyright 2026 by World Scientific Publishing Co.
%% All rights are reserved.
%%%%%%%%%%%%%%%%%%%%%%%%%%%%%%%%%%%%%%%%%%%%%%%%%%%%%%%%%%%%%%%%%%%%%%%%%%%%
%
%\documentclass[wsdraft]{ws-jbcb}
\documentclass{ws-jbcb}

\usepackage[super]{cite}
\usepackage{xcolor}
\usepackage[verbose,hypertexnames=false]{hyperref}
\hypersetup{colorlinks=false,allbordercolors=blue,pdfborderstyle={/S/U/W 1}}
% \label, \ref and \cite commands are highly recommended

\begin{document}

\markboth{F. Author \& S. Author (authors' names)}{Instructions for typing manuscripts (paper's title)}

%%%%%%%%%%%%%%%%%%%%% Publisher's Area please ignore %%%%%%%%%%%%%%%
%
\catchline{}{}{}{}{}
%
%%%%%%%%%%%%%%%%%%%%%%%%%%%%%%%%%%%%%%%%%%%%%%%%%%%%%%%%%%%%%%%%%%%%

\title{Instructions for typesetting manuscripts\\
using \LaTeX\footnote{For the title, try not to use more than
3 lines. Typeset the title in 10~pt Roman, sentence case and
boldface.}
}

\author{First Author\footnote{Typeset names in 8~pt Roman, upper and lower case. Use the footnote to indicate the present or permanent address of the author.}}

\address{University Department, University Name, Address\\
City, State ZIP/Zone,
Country\footnote{State completely without abbreviations, the
affiliation and mailing address, including country. Typeset in 8~pt
Times italic.}\\
first\_author@university.edu}

\author{Second Author\footnote{Use the footnote to indicate the Corresponding author.}}

\address{Group, Laboratory, Address\\
City, State ZIP/Zone, Country\\
second\_author@group.com}

\maketitle

\begin{history}
\received{(Day Month Year)}
\revised{(Day Month Year)}
\accepted{(Day Month Year)}
\published{(Day Month Year)}
\end{history}

\begin{abstract}
The abstract should summarize the context, content
and conclusions of the paper in less than 200 words. It should
not contain any references or displayed equations. Typeset the
abstract in 8~pt Times roman with baselineskip of 10~pt, making
an indentation of 1.5 pica on the left and right margins.
\end{abstract}

\keywords{Keyword1; keyword2; keyword3.}

\section{General Appearance}
Contributions to {\it Journal of Bioinformatics and Computational
Biology} are to be\break in American English. Authors are encouraged
to have their contribution checked\break for grammar. American spelling
should be used. Abbreviations are allowed but should be spelled out in
full when first used. Integers ten and below are to be spelled out.
Italicize foreign language phrases (e.g.~Latin, French).  Upon
acceptance, authors are required to submit their data source file
including postscript files for figures.

The text is to be typeset in 10~pt Times roman, single spaced with
baselineskip of 13~pt. Text area (including copyright block) is 8~inches high and 5~inches wide for the first page.  Text area (excluding running title) is 7.7~inches high and 5~inches wide for subsequent pages. Final pagination and insertion of running titles will be done by the publisher.

\section{Major Headings}
Major headings should be typeset in boldface with the first letter of
important words capitalized.

\subsection{Sub-headings}
Sub-headings should be typeset in boldface italic and capitalize
the first letter of the first word only. Section number to be in
boldface roman.

\subsubsection{Sub-subheadings}
Typeset sub-subheadings in medium face italic and capitalize the
first letter of the first word only. Section numbers to be in
roman.

\subsection{Numbering and spacing}
Sections, sub-sections and sub-subsections are numbered in
Arabic. Use double spacing before all section headings, and
single spacing after section headings. Flush left all paragraphs
that follow after section headings.

\subsection{Lists of items}
Lists may be laid out with each item marked by a bullet:
\begin{itemlist}
 \item item one,
 \item item two,
 \item item three.
\end{itemlist}
Items may also be numbered in lowercase roman numerals:
\begin{romanlist}[(iii)]
 \item item one
 \item item two
	\begin{alphlist}[(b)]
	\item Lists within lists can be numbered with lowercase
              roman letters,
	\item second item.
	\end{alphlist}
\item item three
\end{romanlist}

\section{Equations}
Displayed equations should be numbered consecutively in each
section, with the number set flush right and enclosed in
parentheses
\begin{equation}
\mu(n, t) = {\sum^\infty_{i=1} 1(d_i < t, N(d_i) = n) \over
\int^t_{\sigma=0} 1(N(\sigma) = n)d\sigma}\,.
\label{this}
\end{equation}

Equations should be referred to in abbreviated form,
e.g.~``Eq.~(\ref{this})''. In multiple-line
equations, the number should be given on the last line.

For multiline equations, prefer using \verb|align| or \verb|align*| instead of \verb|eqnarray|.
\begin{align}
g(x) & = \sum_{k=1}^{N} y_k g(x_k) \nonumber\\
& = \sum_{k=1}^{N} y_k \sum_{l=1}^{M} b_{kl} w_l \nonumber\\
& = \sum_{k=1}^{N} \sum_{l=1}^{M} b_{kl} y_k w_l
\end{align}

Displayed equations are to be centered on the page width.
Standard English letters like x are to appear as $x$
(italicized) in the text if they are used as mathematical
symbols. Punctuation marks are used at the end of equations as
if they appeared directly in the text.

\section{Theorem Environments}

\begin{theorem}
\label{thm1}
Theorems$,$ lemmas$,$ etc. are to be numbered consecutively in the
paper. Use double spacing before and after theorems$,$ lemmas$,$ etc.
\end{theorem}

The citation command can be used to cross-link for theorem
declaration, see Theorem~\ref{thm1}.

\begin{proof}
Proofs should end with a box.
\end{proof}

\vspace{-16pt}
\section{Illustrations and Photographs}
Figures are to be inserted in the text nearest their first reference.
Where possible, figures should be embedded in electronic form from
scanned images or a computer-based drawing package. If the author requires the
publisher to reduce the figures, ensure that the figures (including
letterings and numbers) are large enough to be clearly seen after
reduction. If photographs are to be used, only black and white ones
are acceptable.

\begin{figure}[t]
\centerline{\includegraphics[width=5.5cm]{jbcbf1}}
\vspace*{8pt}
\caption{A schematic illustration of dissociative recombination. The
direct mechanism, 4m$^2_\pi$ is initiated when the
molecular ion S$_{\rm L}$ captures an electron with kinetic energy.}
\label{fig1}
\alttext{Description of the figure}
\end{figure}

Figures are to be sequentially numbered in Arabic numerals. The
caption must be placed below the figure. Typeset in 8~pt Times roman
with baselineskip of 10~pt. Use double spacing between a caption and
the text that follows immediately.

\pagebreak

Previously published material must be accompanied by written
permission from the author and publisher.

\subsection{Color illustrations and images}
Please prepare all line drawings, halftones (gray scale) and color
illustrations in high resolution.

\subsubsection{The requirements}
\begin{enumerate}
\item 600 dpi for line drawings (black and white).
\item 300 dpi for halftones (gray scale). Do not convert from color
images as they reproduce very poorly.
\item 300 dpi for color images. Must be in
CMYK (Cyan, Magenta, Yellow and Black) for color separation.
RGB (Red, Green and Blue) is unacceptable for color separation work.
Color images to be printed in color are allowed only as agreed
in the contract.
\item Ensure all labels/annotations are sharp and clear for reproduction.
Easy-to-read typeface/font like Arial, Helvetica or Times Roman are
recommended for the labels.
\item To provide softcopy of the illustrations in either eps, ps, tif,
jpg, gif or bmp format, preferable on a PC platform. All illustrations
should be embedded in the text. At the same time provide the illustration,
example, John$\_$fig1.eps, Hogn$\_$fig2.eps, or
John$\_$fig1.tif, Hogn$\_$fig2.tif.
\end{enumerate}

\subsubsection{Color charge}
Contributors have to pay the following additional cost:
\begin{romanlist}
\item US\$200 for each color illustration;
\item an additional US\$100 for Color Supporting Information that
must be scanned.
\end{romanlist}

A single figure is defined as original art that can be processed as a
unit and printed on one page without intervening type.

\section{Recommended Article Length}

\begin{center}
\vspace{5pt}
\begin{tabular}{@{}l@{\hspace{0.5in}}c@{}}
\Hline\\[-8pt]
{\bf Category} &{\bf Free Reformatted Page} \\[3pt]
\hline\\[-8pt]
Research Paper & 15\\
Critical Comments & 10\\
Survey/Reviews & 15\\
Tutorial Papers & 30\\
In silico studies (open access)	&\phantom{0}0\\[3pt]
\Hline
\end{tabular}
\vspace{10pt}
\end{center}

Please note that {\bf US\$40 per reformatted page} will be charged for
any article exceeding the recommended length.

Articles in ``In silico studies'' category will only be considered for open access publication. Article Processing Charge is charged to make an article in open access publication.

\section{Billing Address}
Billing address and contact numbers, together with the manuscripts, should
be provided during submission. Please indicate illustrations that
are designated for color print.

\begin{table}[b]
\tbl{Comparison of acoustic for frequencies for piston-cylinder problem.}
{\begin{tabular}{@{}cccc@{}}
\toprule
Piston mass & Analytical frequency & TRIA6-$S_1$ model &\% Error \\
& (Rad/s) & (Rad/s) \\
\colrule
1.0\hphantom{00} & \hphantom{0}281.0 & \hphantom{0}280.81 & 0.07 \\[2pt]
0.1\hphantom{00} & \hphantom{0}876.0 & \hphantom{0}875.74 & 0.03 \\[2pt]
0.01\hphantom{0} & 2441.0 & 2441.0\hphantom{0} & 0.0\hphantom{0} \\[2pt]
0.001 & 4130.0 & 4129.3\hphantom{0} & 0.16\\
\botrule
\end{tabular}}
\begin{tabnote}
Table notes
\end{tabnote}
\begin{tabfootnote}
\tabmark{a} Table footnote A\\
\tabmark{b} Table footnote B
\end{tabfootnote}
\end{table}

\section{Tables}
Tables should be inserted in the text as close to the point of
reference as possible. Some space should be left above and below
the table.

Tables should be numbered sequentially in the text in Arabic
numerals. Captions are to be centralized above the tables.
Typeset tables and captions in 8~pt Times roman with
baselineskip of 10~pt.

If tables need to extend over to a second page, the continuation of
the table should be preceded by a caption, e.g.~``{\it Table~1}
$(${\it Continued}$)$''.

\section{References}
References in the text are to be numbered consecutively in
Arabic numerals, in the order of first appearance. They are to
be typed in superscripts after punctuation marks,
e.g.~``$\ldots$ in the statement.\cite{joliat}''

References are to be listed in the order cited in the text. Use
the style shown in the following examples. For journal names,
use the standard abbreviations. Typeset references in 9~pt Times
roman.

\section{Footnotes}
Footnotes should be numbered sequentially in superscript
lowercase roman letters.\footnote{Footnotes should be
typeset in 8~pt Times roman at the bottom of the page.}

\section*{Note Added}
Additional notes can be added before Acknowledgment.

\section*{Acknowledgments}
This section should come before the Appendices. Funding
information may also be included here.

\section*{ORCID}
You are encouraged to include in your user information the ORCID (\url{https://orcid.org/}) or register for one if you don't have it. This ID will help to identify you in the researcher community and make it easier to keep track of all your publications.
Please provide a valid ORCID here, e.g.,

\noindent Josiah Carberry - \url{https://orcid.org/0000-0002-1825-0097}

\noindent Rajesh Babu - \url{https://orcid.org/0009-0006-0415-6880}

\appendix

\section{Appendices}
Appendices should be used only when absolutely necessary. They
should come before the References. If there is more than one
appendix, number them alphabetically. Equations number
occurring in the Appendix is a continuation from the last
equation number of the main text
\begin{equation}
\mu(n, t) = {\sum^\infty_{i=1} 1(d_i < t, N(d_i) = n) \over
\int^t_{\sigma=0} 1(N(\sigma) = n)d\sigma}\,.
\label{appeqn}
\end{equation}

\begin{thebibliography}{0}
\bibitem{beeson} Beeson MJ, {\it Foundations of Constructive Mathematics},
Springer, Berlin, 1985, \url{https://doi.org/10.1007/978-3-642-68952-9}.

\bibitem{clark} Clark KL, Negations as failure, in Gallaire H,
Winker J (eds.), {\it Logic and Data Bases}, Plenum Press, New York,
pp.~293--306, 1973, \url{https://doi.org/10.1007/978-1-4684-3384-5_11}.

\bibitem{tamassia} Tamassia R, Batini C, Talamo M,
An algorithm for automatic layout of entity relationship diagrams,
in Davis CG, Jajodia S, Ng PA, Yeh RT (eds.),
{\it Entity-Relationship Approach to Software Engineering,
Proc. 3rd Int. Conf. Entity-Relationship Approach},
North-Holland, Amsterdam, pp.~421--439, 1983.

\bibitem{gewirtz} Gewirtz WL, Investigations in the theory of
descriptive complexity, Ph. D. Thesis, New York University, 1974.

\bibitem{joliat} Joliat M, A simple technique for partial elimination of
unit productions from LR({\it k}) parsers, {\it IEEE Trans Comput}
{\bf 27}:753--764, 1976, \url{https://doi.org/10.1109/TC.1976.1674686}.

\bibitem{lorentz} Lorentz R, Benson DB,
Deterministic and nondeterministic flow-chart interpretations,
{\it J Comput System Sci} {\bf 27}:400--433, 1983 \url{https://doi.org/10.1016/0022-0000(83)90050-8}.
\end{thebibliography}

\bigskip\bigskip
\noindent
\parbox{5truein}{
\begin{minipage}[b]{1truein}
\centerline{\includegraphics[width=1in,height=1.25in]{john}}
\end{minipage}
\hfill         %to 2nd column
\begin{minipage}[b]{3.85truein}
{\bf Andr\'e Kaup} received the Diploma and Doctoral degrees,
both in electrical engineering, from Aachen University of Technology
(RWTH), Aachen, Germany, in 1989 and 1995, respectively.

\hglue15pt {From 1989 to 1995, he was with the Institute for Communication
Engineering, RWTH, where he was responsible for industrial as well as
academic research projects in the area of high-resolution digital
color imaging, object-based image analysis,\hfilneg}
\end{minipage} } %close for parbox

\vspace*{4.8pt}
\noindent
and coding and models for human perception. In 1995, he joined the
Networks and Video Communications Department, Siemens Corporate
Technology, Munich,\break
Germany, where he is responsible for several
European research projects in the area of very low bit rate video
coding and mobile multimedia communications.

\indent
Dr. Kaup is a member of the Informationstechnische Gesellschaft
(ITG/VDE). He is an active member in ITU-T SG 16 standardization and
Head of the German ISO/MPEG delegation.

\newpage

%\vspace*{13pt}	%SECOND PHOTO
\noindent
\parbox{5truein}{
\begin{minipage}[t]{1truein}
\centerline{\includegraphics[width=1in,height=1.28in]{yaming}}
\end{minipage}
\hfill         %to 2nd column
\begin{minipage}[b]{3.85truein}
{{\bf Yaming Wang} received his M.Sc. degree in Electrical Engineering
from Zhejiang University of Technology, China, and his Ph.D. degree in
Biomedical Engineering from Zhejiang University, China, in 1996 and
2000, respectively.

\hglue 15pt From August 2000 to August 2001, he was at The Hong Kong
University of Science and Technology as a postdoctoral fellow. He is
currently in the Department of Electronic Engineering, Zhejiang
Institute of Science and Technology.}
\end{minipage} } %close for parbox
\vfill\eject

\end{document} 